جاءت توقعات \textbf{كافة الشركات} حول \textbf{مؤشرات الإنتاج والمبيعات المحلية واستغلال الطاقة الإنتاجية} خلال الربع القادم (يناير - مارس 2026) بقيم أعلى \textbf{من المستوى المحايد}، وأقل من الربع \textbf{السابق} وأفضل من الربع \textbf{المناظر}. أما \textbf{الشركات الصغيرة والمتوسطة} فحققت في \textbf{مؤشري الصادرات والمخزون السلعي} توقعات أفضل من الربعين \textbf{السابق والمناظر}. وفي المقابل توقعت \textbf{الشركات الكبيرة} في مؤشر \textbf{الصادرات} بقيم أعلى من \textbf{المستوى المحايد} بـ 17 نقطة، وأقل من الربع \textbf{السابق} بـ 4 نقاط، وأفضل من الربع \textbf{المناظر} بـ 7 نقاط. وسجلت التوقعات لـ\textbf{المخزون السلعي} نفس قيم الربع السابق متجاوزة \textbf{المستوى المحايد} بـ 25 نقطة، والربع \textbf{المناظر} بـ 15 نقطة (الشكل 4-3).
